\chapter{Platform Architecture}
\label{ch:platform}

The research of \cref{ch:equity,ch:energy} runs on a data-engineering
platform that was prior work to this project; the project's architectural
contribution is the pair of seams that let relational factors join that
platform without forking it, and the loop that carries research outputs
back into the warehouse. This chapter describes the platform in the
depth the later chapters rely on --- enough to make every named
component findable in the repository --- and no further.

\section{The layered platform}
\label{sec:pl-layers}

The platform follows a conventional analytics layering. Ingestion
fetches daily equity bars (Yahoo Finance), deterministic synthetic
prices for offline work, and the ENTSO-E power-market series used in
\cref{ch:energy}: day-ahead prices, load and renewable forecasts,
realised load, generation mix, and physical cross-border flows with
transfer capacities \cite{entsoe-tp}. Raw and processed data land as
Parquet files in a local data lake. The warehouse is DuckDB
\cite{raasveldt2019duckdb}, chosen for zero-infrastructure local
reproducibility --- anyone with the repository can rebuild every table.
Transformations are split by nature: factor mathematics lives in Python,
while warehouse marts are dbt models \cite{dbt} so that derived tables
are declarative, tested, and documented. A Kestra flow orchestrates the
daily equity and hourly energy pipelines, and a Streamlit dashboard is
the read-only application layer.

One boundary is stated in the architecture documentation and repeated
here: the platform is a research and portfolio-demonstration system. It
places no orders, manages no broker state, and offers no investment
advice.

\section{Two seams instead of a fork}
\label{sec:pl-seams}

The capstone's design problem was to add a second \emph{family} of
factors --- relational ones --- to a platform whose registries, apply
step, backtests, and marts all assumed island factors. The naive
solution is a parallel pipeline; the project's solution is two narrow
interfaces, fixed in architecture decision records before implementation.

The first seam is the propagator: one snapshot's node features and
topology in, one factor value per node out. Both the unlearned baseline
(uniform neighbour averaging) and the trained GAT are adapters on this
single interface. The seam is deliberately per-snapshot and
transform-only --- training happens elsewhere, weights are bound at
construction --- which keeps it a pure, testable function and keeps
attention behind the interface rather than in it. The design consequence
carried through both result chapters: because baseline and GAT share one
seam, the attention A/B of \cref{sec:anchors} is a one-line adapter
swap, not a comparison across two code paths that might differ in
untracked ways. A characterisation test pins the equivalence: the
uniform propagator on a fully-connected topology must reproduce the
platform's pre-existing cross-market mean exactly.

The second seam unifies factor declaration. The two tracks had grown two
factor idioms --- equity as metadata plus a compute callable, energy as
metadata interpreted by a monolithic imperative function. One
\texttt{Factor} contract over the canonical \emph{(time, entity)} panel,
plus a \texttt{FactorProvider} interface for anything that yields
factors, absorbed both: the energy legacy path is wrapped as one
provider, and the graph factor is simply another, whose compute closes
over a propagator and a topology source and runs the per-snapshot loop
internally. At the registry and apply layer, island and relational
factors are indistinguishable --- which is precisely what makes the
four-gate comparison of \cref{sec:gates} a fair one, since every factor
reaches the evaluation through the same code.

\section{From research run to warehouse mart}
\label{sec:pl-loop}

Research results do not live in notebooks. A persistence step flattens a
GAT run into four warehouse tables --- the factor panel (composite,
anchors, and island alphas per date and symbol), per-alpha diagnostics,
the four gates as a one-row report, and the attention A/B as another ---
and dbt models build the presentation layer on top: a staging view over
the panel, a tiered factor-versus-baseline mart, and a one-row scorecard
joining gates to the A/B. The dashboard reads the marts, not the
research code. The practical effect showed up during this report's own
preparation: rebuilding the marts end-to-end (21 dbt models and tests on
the equity warehouse, 12 on the energy warehouse) is a single command
per project, and the full test suite --- 115 tests, of which 55 belong
to the capstone modules --- is the reproducibility contract for
everything the previous chapters claimed. The leakage-critical pieces
(labels, splits, rank-IC) are deliberately pure pandas so they are
testable without the GNN dependency stack.
